\section{Test problems and exact references}
\label{sec:testcases}

\fig{hero} shows the geodesics of the study. Each test case fixes the initial data, the span
of the integration, the events that stop it, an exact reference and an error defined on a
physical quantity, never on the raw state, because the states of (a) and (b) differ.

\begin{figure}[t]
\centering
\includegraphics[width=\textwidth]{F01_hero.pdf}
\caption{\label{fig:hero}The geodesics of the study in the equatorial plane, with the horizon
($r = 2$), the photon sphere ($r = 3$) and the ISCO ($r = 6$). Left: a parallel bundle of light
rays with impact parameters from $2b_c$ to $b_c(1 + 10^{-6})$; the last winds about twice
around the photon sphere. Right: six radial periods of the bound orbit $(p, e) = (12, 0.5)$.}
\end{figure}

\paragraph{TC1, light deflection.} A photon with impact parameter $b > b_c = 3\sqrt3$ starts at
$r_\mathrm{far}$ and is followed to $r_\mathrm{far}$ again. With $u = 1/r$ and the roots
$u_1 < 0 < u_2 < u_3$ of $2u^3 - u^2 + b^{-2}$, the angle swept is
\begin{equation}
  \Delta\phi(r_\mathrm{far}) = \frac{4F(\psi\,|\,m)}{\sqrt{2(u_3 - u_1)}}, \qquad
  m = \frac{u_2 - u_1}{u_3 - u_1}, \qquad
  \sin^2\psi = \frac{(u_3 - u_1)(u_2 - u_\mathrm{far})}{(u_2 - u_1)(u_3 - u_\mathrm{far})},
  \label{eq:deflection}
\end{equation}
which for $r_\mathrm{far} \to \infty$ is Darwin's result~\cite{darwin1959,iyer2007}. The
reference uses the same $r_\mathrm{far}$ as the integration: comparing with the asymptotic
angle would leave a systematic error of order $b/r_\mathrm{far}^2$. Equation
(\ref{eq:deflection}), Darwin's form in terms of the periapsis and a direct 30-digit quadrature
agree to double precision for $b \in [5.2, 10^4]$. TC1b takes $b = b_c(1 + 10^{-k})$,
$k = 1, \dots, 10$, where the number of windings grows like $-\ln(b/b_c - 1)/2\pi$
(Bozza's strong-deflection limit~\cite{bozza2002}).

\paragraph{TC2, photon sphere.} A photon starts at a turning point $r_0 = 3 + \delta_0$ with
$b = b_c$ and is followed until $|r - 3| > 0.1$. Small offsets grow like
$\cosh(\lambda_L t)$ with $\lambda_L = 1/(3\sqrt3)$~\cite{cardoso2009}, a factor $e^{2\pi}
\approx 535$ per orbit; the exact exit angle is an elliptic integral from the turning point.
The equatorial orbit with $\delta_0 = 0$ is an exact fixed point of every Runge--Kutta map:
the radial components of the vector field vanish and $t$, $\phi$ are linear in $\lambda$. It
survived 900 orbits in a first test. The experiments therefore tilt the orbital plane by
$45^\circ$, so that truncation errors in $L^2$ reach $r$.

\paragraph{TC3, circular orbits.} Timelike circular orbits with $E = f/\sqrt{1 - 3/r_c}$,
$L = \sqrt{r_c/(1 - 3/r_c)}$ and $\Omega = r_c^{-3/2}$. Like the photon sphere they are fixed
points of the radial dynamics; their errors measure round-off, not truncation. TC3m starts on
the ISCO ($r = 6$, $E = \sqrt{8/9}$, $L = 2\sqrt3$) in a plane inclined by $45^\circ$ and stops
when $|r - 6| > 1$.

\paragraph{TC4 and TC5, eccentric orbits.} Bound orbits are labelled by the semi-latus rectum
$p$ and eccentricity $e$~\cite{cutler1994}, with $r_p = p/(1 + e)$ and $r_a = p/(1 - e)$. The
orbit equation $(\dd u/\dd\phi)^2 = 2(u - u_1)(u - u_2)(u - u_3)$ has $u_1 = (1 - e)/p$,
$u_2 = (1 + e)/p$ and $u_3 = \tfrac12 - 2/p$, which gives the angle per radial period and the
orbit in closed form~\cite{chandrasekhar1983}:
\begin{equation}
  \Phi = \frac{4K(m)}{\sqrt{2(u_3 - u_1)}}, \qquad
  u(\phi) = u_1 + (u_2 - u_1)\,\mathrm{cd}^2\!\left(\sqrt{\tfrac{u_3 - u_1}{2}}\,\phi\;\middle|\;m\right),
  \qquad m = \frac{u_2 - u_1}{u_3 - u_1}.
  \label{eq:orbit}
\end{equation}
The radial periods in proper and coordinate time come from a 30-digit quadrature in the angle
$\chi$ of $r = \tfrac12(r_a + r_p) - \tfrac12(r_a - r_p)\cos\chi$, which removes both
turning-point singularities and gives $\Phi$ a second time (agreement $10^{-16}$). The
long-term case TC5 follows $(20, 0.5)$ for $10^3$--$10^4$ periods; its phase error is
$\phi_N - N\Phi$ at the $N$-th periapsis, which needs no reference integration.

\paragraph{TC6, inclined orbits.} The orbit $(20, 0.5)$ in planes inclined by $30^\circ$,
$60^\circ$ and $85^\circ$. The references are $L^2$ and the orientation of the orbital plane,
measured by the angle between the angular-momentum vector $(L_x, L_y, L_z)$ at each sample and
at the start.

\paragraph{Validation.} \tab{validation} compares tightly integrated values (DOP853, tolerance
$10^{-13}$, formulation (b)) with every reference. All agree to $10^{-11}$ or better, except
the cases amplified by the photon sphere, which agree to about $10^{-8}$.

\begin{table*}[t]
\caption{\label{tab:validation}Validation (T2): DOP853 at tolerance $10^{-13}$ in formulation
(b) against the exact references. The last row compares with the value listed by the
predecessor project, \texttt{schwarzschild-geodesics}: its ``analytic'' 2.127093985 was
evaluated at the smallest sampled radius rather than the true periapsis. The exact value,
from the closed form and a 40-digit quadrature, is 2.12709400813.}
\centering\footnotesize
% Generated by scripts/make_tables.py from results/summary/. Do not edit.
\begin{tabular}{lrrr}
\toprule
Check & Numerical & Analytical & Difference \\
\midrule
TC1 deflection $\Delta$$\phi$(r\_far), b = 6, r\_far = 1000 & 4.84898 & 4.84898 & $6.1\times10^{-13}$ \\
TC1 deflection $\Delta$$\phi$(r\_far), b = 40, r\_far = 10000 & 3.2417 & 3.2417 & $4.1\times10^{-13}$ \\
TC1 deflection $\Delta$$\phi$(r\_far), b = 1000, r\_far = 10000 & 2.94527 & 2.94527 & $3.5\times10^{-14}$ \\
TC1b windings, b = b\_c(1 + 1e-6) & 2.63346 & 2.63346 & $1.1\times10^{-8}$ \\
TC2 orbits before $|$r - 3$|$ $>$ 0.1, r0 = 3 + 1e-6 & 1.93916 & 1.93916 & $1.2\times10^{-8}$ \\
TC2 growth rate of $|$r - 3$|$ in t (Lyapunov exponent) & 0.192515 & 0.19245 & $6.5\times10^{-5}$ \\
TC3 circular orbit r\_c = 10: d$\phi$/dt & 0.0316228 & 0.0316228 & $0$ \\
TC3 circular orbit r\_c = 6: d$\phi$/dt & 0.0680414 & 0.0680414 & $2.6\times10^{-15}$ \\
TC4 periapsis advance, (p, e) = (20, 0.5) & 1.23386 & 1.23386 & $5.6\times10^{-13}$ \\
TC4 radial period in proper time, (p, e) = (20, 0.5) & 930.547 & 930.547 & $8.0\times10^{-11}$ \\
TC4 periapsis advance, (p, e) = (7.5, 0.5) & 9.35519 & 9.35519 & $1.3\times10^{-11}$ \\
TC4 radial period in proper time, (p, e) = (7.5, 0.5) & 307.14 & 307.14 & $5.0\times10^{-11}$ \\
TC6 inclined 60$^\circ$: in-plane angle after 2 periods & 15.0341 & 15.0341 & $1.6\times10^{-11}$ \\
TC6 inclined 60$^\circ$: tilt of the orbital plane (rad) & 1.58374e-13 & 0 & $1.6\times10^{-13}$ \\
TC6 inclined 60$^\circ$: max relative change of L$^2$ & 1.35381e-13 & 0 & $1.4\times10^{-13}$ \\
Old README: precession r\_apo = 20, L = 4.2 & 2.12709 & 2.12709 & $2.3\times10^{-8}$ \\
\bottomrule
\end{tabular}

\end{table*}
